% =====================================================================
% 10. Physical climate risk: storm overlay -- question (xi)
% Numbers \val{q11...} come from scripts 06 (CLIMADA) and 07 (overlay), \val{q11x...} from script 07b
% (both credit channels, IRB, hazard stress).
% Relative changes are in parts per million (ppm, 1e-6): basis points are too coarse.
% =====================================================================
\section{Physical Climate Risk: Storm Overlay}
\label{sec:climate}

\paragraph{What kind of climate risk, and what kind of statement.}
Climate risks are usually split into transition risks, which arise from the move to a low-carbon economy (regulation, technology, preferences), and physical risks, which are either acute (abrupt extreme events) or chronic (persistent shifts such as warming or sea-level rise) \citep{Acharya23}. Winter storms are an \emph{acute physical} risk, and the overlay transmits them to credit risk through the collateral, i.e.\ through the loss given default. Transition risk, chronic risk and the other channels through which climate reaches a bank (market, liquidity and operational risk) are outside the scope of the question. The nature of the result also matters. Supervisory climate stress tests report losses \emph{conditional} on a scenario, such as the NGFS pathways \citep{NGFS24}, whose probability is not specified. Our catalogue is instead a probabilistic description of the \emph{present} climate: historical storms and perturbed versions of them, each with an annual frequency. The overlay therefore yields unconditional VaR and ES, like the rest of this paper, and not a scenario loss. Section~\ref{subsec:climstress} adds a stylised high-hazard stress to show how a scenario would change the picture.

\subsection{FRAMEWORK: HAZARD, EXPOSURE, VULNERABILITY}
CLIMADA \citep{AznarSiguan19} computes the impact of a natural hazard from three components, which we set up following \citet{HauserEtAl26}.
\begin{description}
\item[Hazard] A catalogue of events $e$ with annual frequencies $f_e$ and intensity footprints. We use the probabilistic set of European winter storms restricted to Switzerland that CLIMADA distributes through its data interface (\texttt{storm\_europe}, country \texttt{CHE}); the intensity of an event at a location is the maximum gust speed in m/s. The catalogue contains \val{q11.hazard.n_events|int} events on \val{q11.hazard.meta.n_centroids|int} grid points.
\item[Exposure] The assets at risk are the collateral that backs recoveries. Each counterparty is placed at the coordinates given in the portfolio file, and the value exposed is its recovery value $E_kR_k$ (USD~\val{q5.recovery_value|usd} in total).
\item[Vulnerability] An impact function maps intensity into the mean damage ratio, $\MDR=\mathrm{MDD}\times\mathrm{PAA}$: the mean damage degree of the affected assets times the share of assets affected. We use the function calibrated by \citet{Welker21} for building damage from winter storms in Switzerland, which is the one shipped with CLIMADA for this hazard.
\end{description}
$\MDR_k(e)$ is the impact function evaluated at the gust speed of event $e$ at the grid point nearest to counterparty $k$ (at most \val{q11.hazard.nearest_km_max|f1}~km away). This is what CLIMADA's impact engine computes for an exposure of value one at that point.

We evaluate the hazard at the exact coordinates of each counterparty rather than at the level of its canton: the coordinates are the most precise information on where the collateral is, and a canton such as Bern or Grisons spans lowland and alpine terrain with very different gust speeds. An event ``affects'' a counterparty when it causes a positive damage ratio at its location. As a robustness check, the damage ratio is also evaluated at the centroid of the canton and shared by all its counterparties (Section~\ref{subsec:climvariants}).

\subsection{INTERPRETATION OF $\phi_k$}
\label{subsec:phi}
The climate-adjusted recovery rate under event $e$ is $R_k(e)=R_k\,(1-\phi_k\,\MDR_k(e))$.
\begin{itemize}
\item $\phi_k=0$: none of the recovery value is physically located at the exposed site. Recovery comes entirely from assets that a storm cannot damage, such as financial collateral, guarantees or assets held elsewhere, so $R_k(e)=R_k$ for every event and the counterparty is immune to the overlay.
\item $\phi_k=1$: the entire recovery value consists of physical assets at the assigned location. Every unit of damage reduces recovery proportionally, $R_k(e)=R_k(1-\MDR_k(e))$, and an event that destroys the site ($\MDR=1$) leaves nothing to recover.
\end{itemize}
The value $\phi_k=0.5$ used here assumes that half of the recovery value is physically exposed.

Two features of the real world, both stressed in the course as open problems of climate stress testing, argue for an \emph{effective} $\phi_k$ below the nominal share of physical collateral. The first is insurance, which decides who bears the damage. In Switzerland, insurance of buildings against natural hazards, storms included, is mandatory for owners, and in 19 cantons it must be bought from the cantonal building insurer \citep{BAFU20}. If the collateral is insured, its storm damage is largely paid by the insurer and does not reduce the bank's recovery; $\phi_k$ should then measure only the uninsured part of the physically exposed collateral (deductibles, contents, insurance gaps). The second is adaptation. The impact function of \citet{Welker21} is calibrated on the claims of the cantonal building insurer of Zurich, so it already reflects building standards and protection in place today; adaptation over longer horizons would lower it further. Both points make $\phi_k=0.5$ a conservative choice for Swiss real collateral.

\subsection{SIMULATION DESIGN AND WHAT IT IMPLIES}
\label{subsec:climdesign}
The overlay is an independent layer on top of the credit simulation of Sections~\ref{sec:simulation} and \ref{sec:results}.
\begin{enumerate}
\item For each of the $10{,}000$ scenarios and each event $e$ of the catalogue, draw $N_e\sim\mathrm{Poisson}(f_e)$, independently across events and of $(\Theta,\epsilon)$. Event $e$ occurs in that simulated year if $N_e\ge1$.
\item For each counterparty, take the worst event of the year: $m_k=\max\{\MDR_k(e):N_e\ge1\}$, with $m_k=0$ if no event that affects its location occurs.
\item Recompute the loss with the adjusted loss given default, using the \emph{same} default indicators $I_k$ as in the baseline.
\end{enumerate}
With $\phi_k=\phi$ for all counterparties, the loss with overlay is
\begin{equation}
\label{eq:lclim}
L^{\mathrm{clim}}=\sum_{k=1}^{100}E_k\big(1-R_k(1-\phi\, m_k)\big)I_k=L+\phi\,A,\quad A=\sum_{k=1}^{100}E_kR_km_kI_k .
\end{equation}
The overlay never changes who defaults; for each defaulted name it adds the fraction $\phi\,m_k$ of the amount that would otherwise have been recovered. Three results follow before any simulation.

\paragraph{Sign.}
$L^{\mathrm{clim}}\ge L$ in every scenario, so the expected loss, VaR and ES can only increase.

\paragraph{Expected loss.}
Because storms are independent of defaults, $\E[L^{\mathrm{clim}}]-\E[L]=\phi\sum_kE_kR_k\,\pi_k\,\E[m_k]$, and $\E[m_k]$ is available in closed form. Order the events that damage location $k$ by decreasing damage ratio, $m_{k,(1)}\ge m_{k,(2)}\ge\dots$, with frequencies $f_{(1)},f_{(2)},\dots$ The worst event of the year is the $j$-th one if none of the more damaging events occurs and the $j$-th does, so
\begin{equation}
\label{eq:worstlaw}
\P\big[m_k=m_{k,(j)}\big]=e^{-(f_{(1)}+\dots+f_{(j-1)})}\big(1-e^{-f_{(j)}}\big),\quad \P[m_k=0]=e^{-\sum_jf_{(j)}} .
\end{equation}

\paragraph{Tail.}
When damage ratios are small, $\phi A$ is a small perturbation of $L$. The change of a quantile is then the expected perturbation \emph{at} the quantile, and that of ES the expected perturbation \emph{beyond} it \citep{Gourieroux00}, the same derivatives that define Euler risk contributions \citep[Sect.~6.3]{McNeil05}. To first order,
\begin{equation}
\label{eq:firstorder}
\begin{aligned}
\VaR_\alpha(L^{\mathrm{clim}})-\VaR_\alpha(L)&\approx\phi\,\E[A\mid L=\VaR_\alpha(L)],\\
\ES_\alpha(L^{\mathrm{clim}})-\ES_\alpha(L)&\approx\phi\,\E[A\mid L\ge\VaR_\alpha(L)],
\end{aligned}
\end{equation}
and by independence of storms and defaults $\E[A\mid L]=\sum_kE_kR_k\,\E[m_k]\,\P[I_k=1\mid L]$. Two consequences follow. First, under independence the tail is not hit harder than the mean: EL, VaR and ES all rise by a fraction of order $\phi\,\E[m_k]$ of the recovery value of the defaulting names, and differ only through the composition of those names. A rare storm need not coincide with a rare credit event to enter the $99.9\%$ quantile; it enters in proportion to its probability. Second, to first order only the expected damage ratios $\E[m_k]$ matter; their positive dependence across counterparties, since one winter storm covers much of Switzerland, affects VaR and ES only at second order. Both statements rely on the independence between storms and $\Theta$ and are tested by simulation below.

\paragraph{A remark on horizons.}
The credit loss refers to a horizon of one week, whereas $N_e\sim\mathrm{Poisson}(f_e)$ with annual frequencies describes the storms of one year. We combine the two, which amounts to assuming that the damage of a whole storm season is unrepaired when the default occurs. Since lower frequencies can only lower every expected damage ratio, this is the conservative convention; the horizon-consistent alternative, with weekly frequencies $f_e/52$, is reported as a variant.

\subsection{SENSITIVITY TO THE DAMAGE RATIO}
Table~\ref{tab:climsens} first reports VaR and ES when every counterparty suffers the same damage ratio $m$ in every scenario, an upper bound for any catalogue whose damage ratios do not exceed $m$. The effect is almost linear in $m$: a uniform damage ratio of $1\%$ raises VaR by \val{q11s.M2.dVaR_0.01|pct2}\% under M2, \val{q11s.M3.dVaR_0.01|pct2}\% under M3 and \val{q11s.M1.dVaR_0.01|pct2}\% under M1, and one of $10\%$ by ten times as much. By (\ref{eq:firstorder}) the elasticity of ES is $\phi$ times the ratio between the recovery value and the loss of the names that default in the tail, $\sum_kE_kR_kI_k/\sum_kE_k(1-R_k)I_k$, which is \val{q11s.M2.tail_recovery_ratio|f2} in the ten tail scenarios of M2. Since $\phi=0.5$, the rule of thumb is that VaR and ES rise by somewhat less than half of the damage ratio.

\begin{table}[tb]
\caption{Upper-bound sensitivity: increase of the $99.9\%$ VaR and ES (in \%, $N=10{,}000$, thresholds $\dhat$) when all counterparties suffer the same mean damage ratio $m$ in every scenario, $\phi_k=0.5$.}
\label{tab:climsens}
\centering
\qrmtable{tab_climate_sensitivity}
\end{table}

\subsection{THE STORM CATALOGUE}
Table~\ref{tab:climhaz} and Figure~\ref{fig:climhaz} summarise the catalogue at the locations of the counterparties. The hazard set is the WISC probabilistic set at a resolution of 4.4~km: it consists of \val{q11.hazard.meta.n_orig_events|int} historical winter storms and 29 perturbed versions of each, \val{q11.hazard.n_events|int} events in total, each with the same annual frequency $1/\val{q11.hazard.return_period_rarest|int}$. Of these, \val{q11.hazard.n_damaging|int} cause a positive damage ratio for at least one counterparty. Their frequencies sum to \val{q11.hazard.sum_freq_damaging|f2} per year, so a damaging event occurs in \val{q11.hazard.p_year_damage|pct0}\% of the years. Damage ratios, however, are tiny. The largest gust speed at any counterparty in the whole catalogue is \val{q11.hazard.gust_max|f1}~m/s, reached during the storm of \val{q11.hazard.meta.top_event_date|date} (Lothar); the largest damage ratio is \val{q11.hazard.mdr_max|pct2}\%; and the expected worst damage ratio of a year from (\ref{eq:worstlaw}) is \val{q11.hazard.exp_worst_mean|ppm} parts per million (ppm) on average across counterparties, ranging from practically zero for the names in canton \val{q11.hazard.canton_low|str}, south of the Alps, to \val{q11.hazard.exp_worst_max|ppm}~ppm in canton \val{q11.hazard.canton_high|str}, on the Jura.

The reason is the impact function of \citet{Welker21} (Figure~\ref{fig:climhaz}, right). It is zero up to about 20~m/s and rises steeply and convexly above, reaching a damage ratio of $0.1\%$ only close to 40~m/s, whereas the gust speeds at the counterparties lie mostly below 20~m/s (99\% quantile among the damaging events: \val{q11.hazard.gust_p99_damaging|f1}~m/s). Most storms thus cause no damage or damage ratios of the order of $10^{-6}$ to $10^{-5}$, and the hazard of a year is concentrated in a few events. The annual damage to the recovery value of the portfolio, $D=\sum_kE_kR_km_k/\sum_kE_kR_k$, is zero in \val{q11.hazard.D_share_zero|pct0}\% of the simulated years, its median is \val{q11.hazard.D_median_over_mean|pct1}\% of its mean, and the worst $1\%$ of the years account for \val{q11.hazard.D_share_of_mean_top1pct|pct0}\% of the mean. The damage ratios of different counterparties are positively but far from perfectly dependent, with an average pairwise correlation across events of \val{q11.hazard.mdr_pairwise_corr_mean|f2}: the most damaging event affects \val{q11.hazard.share_names_hit_by_top_event|pct0}\% of the counterparties at once but destroys only \val{q11.hazard.worst_event_recovery_share|ppm}~ppm of the recovery value of the portfolio.

\begin{table}[tb]
\caption{The winter-storm catalogue at the locations of the 100 counterparties (CLIMADA \texttt{storm\_europe} hazard for Switzerland, WISC event set, impact function of \citet{Welker21}).}
\label{tab:climhaz}
\centering
\qrmtable{tab_climate_hazard}
\end{table}

\begin{figure}[tb]
\centering
\qrmfigure{fig_climate_hazard}
\caption{Left: gust footprint of the most damaging event of the catalogue (26 December 1999) and location of the counterparties (circle area proportional to exposure). Right: impact function (mean damage ratio, logarithmic scale) and frequency-weighted distribution of the gust speeds at the counterparties.}
\label{fig:climhaz}
\end{figure}

\subsection{RESULTS}
\label{subsec:climresults}
Because the effects are small, Table~\ref{tab:clim} reports them in USD and in parts per million of the value without overlay; one basis point is $100$~ppm.

\begin{table}[tb]
\caption{$99.9\%$ VaR and ES without the storm overlay (USD) and their increase with the overlay, in USD and in parts per million ($\phi_k=0.5$).}
\label{tab:clim}
\centering
\qrmtable{tab_climate}
\end{table}

\paragraph{Expected loss.}
The overlay raises the expected loss from USD~\val{q11.EL_base|usd2} to USD~\val{q11.EL_clim_exact|usd2}, by \val{q11.dEL_exact|ppm}~ppm (exact, from (\ref{eq:worstlaw})).

\paragraph{Baseline run.}
On the $10{,}000$ credit scenarios of the baseline run of Section~\ref{sec:results}, a damaging storm occurs in \val{q11.assign.share_years_with_damage|pct0}\% of the simulated years and a counterparty is affected in \val{q11.assign.share_name_years_hit|pct0}\% of the counterparty-years, with an average worst damage ratio of \val{q11.assign.mean_worst_mdr_if_hit|ppm}~ppm when affected. The overlay raises VaR by USD~\val{q11.assign.M1.dVaR_usd|usd2} under M1, USD~\val{q11.assign.M2.dVaR_usd|usd2} under M2 and USD~\val{q11.assign.M3.dVaR_usd|usd2} under M3, and ES by USD~\val{q11.assign.M1.dES_usd|usd2}, USD~\val{q11.assign.M2.dES_usd|usd2} and USD~\val{q11.assign.M3.dES_usd|usd2} (Table~\ref{tab:clim}). In relative terms no change exceeds \val{q11.assign.M2.dES|ppm}~ppm.

These figures should not be compared across models: with $N=10{,}000$ the VaR is a single scenario, the tenth-largest loss, and its increase is the add-on of the one storm year paired with it, usually close to zero and occasionally large given the skewness of the annual damage. To measure this sampling error we keep the credit scenarios fixed and redraw the $10{,}000$ storm years \val{q11.rep.n_rep|int} times (Figure~\ref{fig:climeff}, top). Under M2 the increase of VaR is zero in \val{q11.rep.M2.dVaR_share_zero|pct0}\% of the re-draws and has a 95\% quantile of \val{q11.rep.M2.dVaR_q95|ppm}~ppm, but its mean is \val{q11.rep.M2.dVaR_mean|ppm}~ppm with a standard deviation of \val{q11.rep.M2.dVaR_sd|ppm}~ppm. The increase of ES, an average over ten scenarios, has a mean of \val{q11.rep.M2.dES_mean|ppm}~ppm and a standard deviation of \val{q11.rep.M2.dES_sd|ppm}~ppm. The other two models behave alike.

\paragraph{Reference run and the first-order effect.}
The first-order approximation (\ref{eq:firstorder}) uses the exact $\E[m_k]$ and needs no storm simulation. It gives an increase of \val{q11.first_order.M1.dVaR|ppm}, \val{q11.first_order.M2.dVaR|ppm} and \val{q11.first_order.M3.dVaR|ppm}~ppm for VaR and \val{q11.first_order.M1.dES|ppm}, \val{q11.first_order.M2.dES|ppm} and \val{q11.first_order.M3.dES|ppm}~ppm for ES under M1, M2 and M3 (first row of Table~\ref{tab:climvar}). It is confirmed by the simulations in two ways. First, it coincides with the \emph{average} increase over the storm re-draws of the baseline run (\val{q11.rep.M2.dVaR_mean|ppm}~ppm for VaR and \val{q11.rep.M2.dES_mean|ppm}~ppm for ES under M2). Second, in the reference run, where the two million credit scenarios are each paired with an independent storm year, the increase of ES is \val{q11.ref.base.M1.dES|ppm}, \val{q11.ref.base.M2.dES|ppm} and \val{q11.ref.base.M3.dES|ppm}~ppm, within about 15\% of the first-order values. The increase of VaR in the reference run, between zero and \val{q11.ref.base.M1.dVaR|ppm}~ppm, is again the add-on of a single scenario: a quantile responds to the expected add-on only on average over storm draws, and for this hazard a single draw is almost always below the average.

Three conclusions follow. \emph{First}, the tail is not hit harder than the mean. The increase of ES is about \val{q11.first_order.M2.dES|ppm}~ppm against \val{q11.dEL_exact|ppm}~ppm for the expected loss. By (\ref{eq:firstorder}) the relative effect is $\phi$ times the expected damage ratio times the ratio between the recovery value and the loss given default of the defaulting names. That ratio is \val{q11.hazard.recovery_ratio_EL_weighted|f2} for the names that generate the expected loss, mostly B and CCC names, and \val{q11s.M2.tail_recovery_ratio|f2} for the names that default in the tail scenarios of M2. The difference in composition accounts for the whole gap, since $\val{q11s.M2.tail_recovery_ratio|f2}/\val{q11.hazard.recovery_ratio_EL_weighted|f2}$ is close to the ratio of the two relative increases. \emph{Second}, the effect is economically negligible. An increase of ES by USD~\val{q11.first_order.M2.addon_tail_usd|usd2} under M2 compares with a Monte Carlo standard error of VaR of USD~\val{q8.M2.se_VaR_10k_q|usd} at $N=10{,}000$ and with a difference of USD~\val{q8.M3_minus_M2.dVaR_ref|usd} between M3 and M2: it is four orders of magnitude below either. \emph{Third}, the catalogue is far from the bound of Table~\ref{tab:climsens}. A uniform damage ratio of $0.1\%$, a third of the largest damage ratio in the catalogue, would raise VaR by \val{q11s.M2.dVaR_0.001|ppm}~ppm under M2. The actual effect is two orders of magnitude smaller, because the largest damage ratios occur in few events and at few locations.

\begin{figure}[tb]
\centering
\qrmfigure{fig_climate_effect}
\caption{Top: increase of the $99.9\%$ VaR (plain) and ES (hatched) due to the overlay with $N=10{,}000$ over \val{q11.rep.n_rep|int} re-draws of the storm years; box: quartiles; whiskers: 5\% and 95\% quantiles; dot: the baseline run; horizontal line: first-order approximation (\ref{eq:firstorder}). Bottom: increase of ES in the reference run for the baseline overlay and the variants of Table~\ref{tab:climvar}.}
\label{fig:climeff}
\end{figure}

\subsection{VARIANTS AND A MORE DIVERSIFIED PORTFOLIO}
\label{subsec:climvariants}
\begin{table}[tb]
\caption{Reference run ($N=2{,}000{,}000$): increase of the expected loss (under M2), VaR and ES due to the storm overlay, in parts per million, for the baseline design and five variants. First row: first-order approximation (\ref{eq:firstorder}) for the baseline. The VaR columns are dominated by the add-on of a single scenario; ES is the reliable measure.}
\label{tab:climvar}
\centering
\qrmtable{tab_climate_variants}
\end{table}

Table~\ref{tab:climvar} and the lower panel of Figure~\ref{fig:climeff} report five variants of the baseline, each changing one assumption. Since the increase of VaR is dominated by the add-on of a single scenario, we compare the variants through the expected loss and ES.

\paragraph{Canton-level damage.}
With the damage ratio evaluated at the centroid of the canton, the expected loss rises by \val{q11.ref.canton.M2.dEL|ppm}~ppm against \val{q11.ref.base.M2.dEL|ppm}~ppm in the baseline, and ES under M2 by \val{q11.ref.canton.M2.dES|ppm}~ppm against \val{q11.ref.base.M2.dES|ppm}~ppm. The correlation between the two sets of damage ratios across events and counterparties is only \val{q11.hazard.corr_point_canton|f2}, because gusts vary strongly over a few kilometres, but the errors average out over the portfolio: the choice between coordinates and cantons does not affect the conclusions.

\paragraph{Collateral sensitivity.}
By (\ref{eq:lclim}) the add-on is linear in $\phi$. With $\phi_k=1$, the case in which all recovery value is physically exposed, the expected loss rises by \val{q11.ref.phi1.M2.dEL|ppm}~ppm and ES under M2 by \val{q11.ref.phi1.M2.dES|ppm}~ppm, twice the baseline.

\paragraph{Horizon.}
With weekly frequencies $f_e/52$ the expected loss rises by \val{q11.dEL_exact_weekly|ppm}~ppm (exact), against \val{q11.dEL_exact|ppm}~ppm with annual frequencies. Lower frequencies can only lower every $\E[m_k]$, so the annual convention is the conservative one.

\paragraph{A more geographically diversified portfolio.}
All counterparties are located in Switzerland, a territory smaller than the footprint of a typical European winter storm, so one event affects much of the portfolio at once. To see what geographical diversification would change, we simulate the limiting case in which every counterparty draws its own, independent storm year from the exact law (\ref{eq:worstlaw}). Each location keeps its hazard, hence its expected damage ratio and the expected loss, but no two names are ever hit by the same event. The first-order analysis predicts little change, because VaR and ES respond to the \emph{expected} damage at the locations of the defaulting names, and the simulation confirms it: the increase of ES is \val{q11.ref.diversified.M1.dES|ppm}, \val{q11.ref.diversified.M2.dES|ppm} and \val{q11.ref.diversified.M3.dES|ppm}~ppm under M1, M2 and M3, against \val{q11.ref.base.M1.dES|ppm}, \val{q11.ref.base.M2.dES|ppm} and \val{q11.ref.base.M3.dES|ppm}~ppm in the baseline, higher for one model and lower for the others, with no systematic difference.

A more diversified portfolio would therefore show a different overlay effect mainly to the extent that the \emph{average} hazard at its locations differs from the Swiss one: smaller in regions with weaker storms, as south of the Alps in our own data, and larger where the hazard is more destructive, such as along the North Sea coasts for winter storms, or in tropical-cyclone and flood zones. It would not be smaller merely because the hazard is spread across regions. Diversification reduces the dispersion of the add-on, which matters for VaR and ES only when damage ratios are large enough to reorder the tail scenarios; for a Swiss winter-storm portfolio they are not.

\paragraph{Dependence between storm years and credit years.}
The conclusion changes if the independence assumption is dropped. In the last variant the storm years are the same as in the baseline but are paired with the credit scenarios in a comonotone way: the most severe storm years (by expected add-on) are assigned to the scenarios with the largest credit losses. The marginal laws of storms and of credit losses are unchanged. The increase of ES under M2 becomes \val{q11.ref.comonotone.M2.dES|ppm}~ppm, or USD~\val{q11.ref.comonotone.M2.dES_usd|usd}, \val{q11.ref.comonotone.M2.dES_vs_base|f0} times the baseline, and that of VaR \val{q11.ref.comonotone.M2.dVaR|ppm}~ppm, while the expected loss rises by \val{q11.ref.comonotone.M2.dEL|ppm}~ppm. Under dependence the tail scenarios receive the damage of the worst storm years and not the average damage. It is in this case that geographical concentration becomes costly: the worst storm year hits all Swiss names together, whereas in a diversified portfolio no single year is the worst everywhere. Even this extreme pairing, however, moves ES by less than $0.02\%$. The pairing is not realistic, since a winter storm does not cause a recession, but it bounds the cost of the independence assumption, the strongest simplification of the overlay, and it is the relevant benchmark for hazards that damage the regional economy on which the obligors depend.

\subsection{LGD AND PD CHANNELS OF PHYSICAL RISK}
\label{subsec:climbis}
\citet{Pozdyshev25} incorporate physical risk into the Vasicek model behind the IRB formula through a binary event: with probability $q$ the asset value of the borrower is scaled by $e^{-\alpha}$. The event raises the loss given default to $\mathrm{LGD}_1=\mathrm{LGD}_0+(1-e^{-\alpha})(1-\mathrm{LGD}_0)$, since it destroys the same fraction of the assets from which the bank recovers, and it raises the default probability to the $q$-deformed mixture $(1-q)\Phi(C)+q\,\Phi(C+\alpha/\sigma_A)$, where $C=\Phi^{-1}(\pi_k)$ and $\sigma_A$ is the asset volatility over the horizon of $\pi_k$. Their IRB capital keeps the Basel form and is multiplied by the LGD add-on $F=1+q(\mathrm{LGD}_1-\mathrm{LGD}_0)/\mathrm{LGD}_0$.

\paragraph{Our overlay is the LGD channel of this model.}
Writing the climate-adjusted loss given default as $1-R_k(e)=\mathrm{LGD}_{0,k}+\phi_k\MDR_k(e)\,(1-\mathrm{LGD}_{0,k})$ shows that the assignment's specification is exactly their LGD channel with $1-e^{-\alpha}=\phi_k\MDR_k(e)$. The catalogue delivers the two inputs that their method takes from external providers: for each counterparty, the probability of at least one damaging storm in a year, $q_k$, between \val{q11x.bis.q_min|pct1}\% and \val{q11x.bis.q_max|pct0}\%, and the jump $\alpha_k=-\log(1-\phi\,\bar m_k)$, with $\bar m_k$ the expected worst damage ratio given damage, at most \val{q11x.bis.alpha_max|sci} in log units. The add-on factors are correspondingly close to one (largest $F_k-1$: \val{q11x.bis.F_max|excess_ppm}~ppm), and the BIS add-on reproduces the exact expected-loss effect of our overlay, $\sum_kE_k\pi_k\mathrm{LGD}_{0,k}(F_k-1)=\phi\sum_kE_kR_k\pi_k\E[m_k]$, to machine precision (relative difference \val{q11x.bis.identity_rel_error|sci}).

\paragraph{The PD channel.}
The assignment keeps default indicators unchanged, but the same storm that damages the collateral also lowers the value of the firm. We add this channel by letting the worst damage of the year move the asset value of counterparty $k$ by $\alpha=-\log(1-\phi\,m_k)$: in the standardised threshold model of Section~\ref{sec:model} the default threshold moves up by $\alpha/\sigma_A$ standard deviations, i.e.\ $I_k=\ind\{Y_k\le d_k+s_k\,\alpha/\sigma_A\}$. We set $\sigma_A=25\%$, the annual asset volatility of the course example, consistent with the one-year $\pi_k$ and the annual storm frequencies, and report $\sigma_A=10\%$ for low-volatility borrowers. With storms independent of $\Theta$, the expected loss with both channels is available in closed form from the law (\ref{eq:worstlaw}). For the tail we use the first-order approach of (\ref{eq:firstorder}): the threshold shifts are of the order of $10^{-4}$ standard deviations, so in the two million reference scenarios they add a default in the tail only a handful of times and the simulated change of ES is pure discreteness noise. We therefore replace the defaults added by the shift with their conditional expectation given $\Theta$, $\sum_kE_k\big(p_k(\Theta;d_k+\text{shift})-p_k(\Theta;d_k)\big)\mathrm{LGD}_k(e)$, while the LGD part uses the simulated defaults. This estimator reproduces the exact expected-loss effect (M2: \val{q11x.ref.pd.M2.dEL_cond|ppm} against \val{q11x.ref.pd.dEL_exact|ppm}~ppm for the PD channel).

\begin{table}[tb]
\caption{Physical risk through the LGD and the PD channel (reference run, $N=2{,}000{,}000$, storms independent of $\Theta$, $\phi_k=0.5$). $\Delta$EL is exact; $\Delta$ES is the first-order effect, with the defaults added by the PD channel replaced by their expectation given $\Theta$. Lower block: every gust speed of the catalogue raised by the stated percentage, both channels, $\sigma_A=25\%$.}
\label{tab:climchannels}
\centering
\qrmtable{tab_climate_channels}
\end{table}

Table~\ref{tab:climchannels} gives three results. \emph{First}, the PD channel is larger than the LGD channel: it accounts for \val{q11x.pd_share_EL|pct0}\% of the expected-loss effect of both channels together (\val{q11x.ref.pd.dEL_exact|ppm} against \val{q11x.ref.lgd.dEL_exact|ppm}~ppm). The reason is an amplification. To first order, the LGD channel raises the expected loss of name $k$ by the relative amount $\phi m_k\,R_k/\mathrm{LGD}_{0,k}$, of order $\phi m_k$. The PD channel shifts the threshold by $\phi m_k/\sigma_A$ standard deviations, and a small default probability responds to a shift with the elasticity $\varphi(C)/\Phi(C)$, so the relative increase is $\phi m_k\,\varphi(C)/(\Phi(C)\,\sigma_A)$: about $11\,\phi m_k$ for $\pi_k=1\%$ and $\sigma_A=25\%$. \emph{Second}, the tail is again hit less than the mean: with both channels ES rises by \val{q11x.ref.both.M2.dES|ppm}~ppm under M2, against \val{q11x.ref.both.dEL_exact|ppm}~ppm for the expected loss, because the extra defaults are spread over all scenarios independently of $\Theta$. \emph{Third}, the total remains economically negligible: USD~\val{q11x.ref.both.M2.dES_usd|usd2} of ES under M2, and \val{q11x.ref.both_s10.M2.dES|ppm}~ppm with $\sigma_A=10\%$.

\paragraph{Climate-adjusted IRB capital.}
Table~\ref{tab:climirb} applies the same parameters to the IRB capital of Section~\ref{sec:irb}, with the stressed default probability computed from the $q$-deformed mixture over the worst damage of the year, the event being independent of the systematic factor. The LGD add-on raises $K$ by \val{q11x.irb.s25.d_lgd|ppm}~ppm, both channels by \val{q11x.irb.s25.d_both|ppm}~ppm (USD~\val{q11x.irb.s25.dRWA_both_usd|usd2} of RWA). In the example of \citet{Pozdyshev25} the mixture model gives much less capital than plugging the climate-adjusted PD into the standard formula, because a large and rare jump is a separate state that is diversified away. Here the opposite occurs: the naive plug-in gives \val{q11x.irb.s25.d_naive|ppm}~ppm, slightly less. With frequent and tiny damages ($q_k$ up to \val{q11x.bis.q_max|pct0}\%, $\alpha_k/\sigma_A$ at most \val{q11x.bis.alphatilde_max|sci}) the mixture is practically a deterministic shift of the threshold, which the Vasicek formula treats like any other increase of the PD, and the plug-in is then smaller only because the regulatory correlation and maturity adjustment decrease with the PD. The diversification benefit of the BIS model requires severe, rare events, which Swiss winter storms are not.

\begin{table}[tb]
\caption{IRB capital of Section~\ref{sec:irb} with physical risk in the formula of \citet{Pozdyshev25}: LGD add-on only, both channels (stressed PD from the $q$-deformed mixture, event independent of the systematic factor), and the climate-adjusted PD plugged into the standard formula.}
\label{tab:climirb}
\centering
\qrmtable{tab_climate_irb}
\end{table}

\subsection{A HIGH-PHYSICAL-RISK STRESS}
\label{subsec:climstress}
A scenario-conditional analysis would replace the present-climate catalogue with one consistent with a warming pathway. The CLIMADA data interface distributes winter-storm sets derived from CMIP6 climate models under SSP2-4.5, SSP3-7.0 and SSP5-8.5; a proper use of them compares the historical and future runs of the same climate model, which is beyond the scope of this question. We instead apply a stylised stress that needs no new hazard data: every gust speed of the catalogue is raised by $5\%$, $10\%$ and $20\%$, deliberately large values that stand in for a hot-house-world outcome, and the damage ratios are recomputed with the same impact function.

\begin{figure}[tb]
\centering
\qrmfigure{fig_climate_stress}
\caption{Left: impact function of \citet{Welker21} (logarithmic scale) and the largest gust speed at any counterparty, as in the catalogue and raised by $5\%$, $10\%$ and $20\%$. Right: increase of the expected loss (exact) and of ES (first order) with both channels, as a function of the gust stress.}
\label{fig:climstress}
\end{figure}

Because the impact function is steep and convex above $20$~m/s, the damage responds far more than proportionally (lower block of Table~\ref{tab:climchannels} and Figure~\ref{fig:climstress}). A $5\%$ increase in gusts multiplies the average expected damage ratio by \val{q11x.stress.s005.Em_ratio_to_base|f1} and a $20\%$ increase by \val{q11x.stress.s020.Em_ratio_to_base|f1}; the largest damage ratio rises from \val{q11x.stress.s000.mdr_max|pct2}\% to \val{q11x.stress.s020.mdr_max|pct1}\%. The expected loss and ES follow almost one for one: under the $20\%$ stress, with both channels, ES rises by \val{q11x.ref.g20.M2.dES|ppm}~ppm under M2, \val{q11x.g20_ES_mult_M2|f0} times the present-climate value. The size of the overlay is thus governed by the hazard scenario far more than by the credit model. Yet even this stress leaves ES within $0.01\%$ of its value, USD~\val{q11x.ref.g20.M2.dES_usd|usd0} under M2: for a Swiss portfolio, the conclusion that winter storms are immaterial for credit risk under independence is robust to the hazard as well.
